\chapter{Glossary}\label{app:glossary}

\begin{description}[style=nextline,leftmargin=0pt,font=\normalfont\bfseries]
\item[FESOM family] The unstructured-mesh ocean models the engine principally targets:
  FESOM2 and FESOM-C. Used here as a collective term; statements specific to one model name it.
\item[FESOM2] A global ocean model on a purely triangular mesh with a fixed depth-level
  (z-level) vertical coordinate. The running example uses a FESOM2 mesh.
\item[FESOM-C] A coastal ocean model on a mixed mesh --- elements may be triangles or
  quadrilaterals --- with a terrain-following (sigma) vertical coordinate.
\item[Node] A vertex of the mesh. Node-layout fields (e.g.\ \code{unod}/\code{vnod},
  scalars, sea-surface height) are defined at nodes.
\item[Element] A mesh face: a triangle (FESOM2) or a triangle or quadrilateral (FESOM-C, HEALPix,
  NEMO). Element-layout velocities are defined at element centroids.
\item[Interface, layer] The vertical level boundaries (interfaces) and the volumes between
  them (layers); $K$ interfaces bound $K-1$ layers. The vertical velocity lives at
  interfaces, the three-dimensional horizontal velocity and scalars at mid-layers.
\item[One-ring] The set of elements that share a given node; the stencil used by the
  localization search and several interpolators.
\item[Barycentric coordinates] The area-ratio weights of a point within a triangle, used to
  blend the three vertex values; the triangular interpolation stencil.
\item[Wachspress coordinates] The generalization of barycentric coordinates to a convex
  polygon, used to interpolate over a quadrilateral element; computed in a division-free product form
  that is exact on edges, which is why quadrilaterals must be convex.
\item[z-level coordinate] A vertical coordinate with fixed interface depths, the same
  everywhere (FESOM2).
\item[Sigma coordinate] A terrain-following vertical coordinate whose interface depths are
  fractions of the local column thickness $H=\max(d+\eta,D_{\min})$ (FESOM-C).
\item[Node-layout / element-layout] Whether a velocity field is stored at nodes or at
  element centroids; node layout admits only the barycentric scheme, element layout all three.
\item[Fused / unfused kernel] The two interchangeable implementations of a step: a single
  compiled loop over the particles, or separate parallel passes; they agree bit-for-bit where both exist.
\item[PCHIP] Monotone piecewise-cubic Hermite reconstruction of the velocity in time, using
  four snapshots; raises the realized temporal order above the linear reconstruction's cap.
\item[Counter-based RNG] A stateless random-number generator in which each draw is a hash of
  an integer key, giving results independent of thread count and reproducible across restarts.
\item[Beaching] The optional stranding of particles in shallow or coastal water, with an
  optional resuspension; the only reversible non-alive status.
\item[\code{time\_dir}] The integration-direction factor ($+1$ forward, $-1$ backward),
  folded into the interpolated velocity so the step size stays positive.
\item[Drifter / particle] A passive Lagrangian tracer advected by the flow; ``drifter'' is
  used in the schematics for the point being interpolated to.
\item[HEALPix] A hierarchical equal-area grid of the sphere on which some ocean products are stored; the
  engine builds a mesh from its pixel centres (dual mesh) or from its pixels (pixel mesh).
\item[Grid hash] The uniform lattice of cells over which mesh nodes are bucketed for the nearest-node
  search; its cell is chosen by measurement for each run.
\item[Exact anchor] A nearest node proven to be the nearest by the cap certificate of the search, the
  default since September 2026.
\item[Edge walk] Tracking the containing element from step to step by crossing element edges, instead of
  searching around the nearest node every time.
\item[Edge least squares] The velocity interpolation that fits a plane to the containing element and its
  edge neighbours, with mirrored ghost values at the coast.
\item[Adjoint process] The backward-in-time counterpart of an advection--diffusion process: the advection
  reversed, the diffusive drift and noise unchanged.
\item[Resumable chain] A long run split into consecutive batch jobs, each resuming from the previous one's
  restart; the concatenated result equals one uninterrupted run.
\item[Shard] One contiguous block of a large release, run as a separate process and merged back afterwards.
\item[Agent skill] A packaged set of instructions, reference files and check scripts that a coding agent
  loads to carry out one kind of task on the engine (\cref{ch:skills}).
\end{description}
